\subsection{Conteos que sustentan la comparación}
\label{sec:conteos_a}
La Tabla~\ref{tab:conteos_a} presenta las obligaciones calculadas con relevo y los conteos de ejecución preservados. La columna Jobs cuenta los jobs científicos. Los valores corresponden a una ejecución y coinciden en las tres repeticiones de cada condición.

\begin{center}
\begin{minipage}{\linewidth}
\centering\small
\captionsetup{hypcap=false}
\captionof{table}{Conteos por workflow en las condiciones seleccionadas}
\label{tab:conteos_a}
\begin{tabular}{@{}lrrrr@{}}
\toprule Caso & Jobs & Obligaciones & Entregas & Direcciones conciliadas\\\midrule
P01 & 1 & 2 & 2 & 2\\
P02 & 1 & 2 & 2 & 2\\
L01 & 3 & 2 & 6 & 6\\
GL2 & 2 & 2 & 5 & 4\\
GL3 & 1 & 2 & 4 & 2\\
I40 & 3 & 2 & 6 & 6\\
A01 & 5 & 5 & 13 & 10\\
A02 & 5 & 13 & 13 & 10\\
D02 & 5 & 13 & 13 & 10\\
GD4 & 2 & 13 & 13 & 4\\
AC1 & 5 & 7 & 10 & 10\\
AC2 & 3 & 7 & 8 & 6\\
AC4 & 2 & 7 & 7 & 4\\
R01-E & 10 & 8 & 26 & 20\\
R02-E & 10 & 26 & 26 & 20\\
\bottomrule
\end{tabular}
\end{minipage}
\end{center}

En A02, D02, GD4, AC4 y R02-E coinciden obligaciones y entregas, además de los bytes. Bajo cobertura inyectiva por identidad y destino, la igualdad de cardinalidad hace biyectiva la correspondencia: no quedan entregas adicionales. La igualdad de conteos o bytes por sí sola no descarta sustituciones. En GL3, dos direcciones contienen cuatro entregas; en GD4, agrupar conserva trece entregas.